\begin{theorem}[Large-order current-rung elementarity]
Let \(G=(V,E)\) be a finite loopless multigraph with its individual edge
identities, let \(e_0\in E\), and let the palette have \(k\) colors.  Let

\[
 \mathcal E:\noderef{ExtendedTashkinovTree}(G,k,e_0,\rho,T,\varphi)
\]

be an arbitrary-input positive-rung ETT with displayed tree \(T\), generating
coloring \(\varphi\), and rung number \(n\).  Assume
\noderef{MaximumProperty} for \(\mathcal E\), a
\noderef{TechnicalInductionContext} through rung \(n\), and
\[
 \noderef{MaximumDegree}(G)+1\leq k,
 \qquad 11\leq \noderef{TashkinovOrder}(G).
\]
Then
\[
 \noderef{ColoredVertexSetElementary}
   \bigl(G,\varphi,\noderef{TreeSequenceVertexSet}(T)\bigr).
\]
In particular, no hierarchy base, good hierarchy, mod-coloring, terminal
closure, or missing-color estimate is assumed.
\end{theorem}
\begin{proof}
Apply \noderef{CurrentElementaryOrGoodHierarchyConstructionLarge} to
\(\mathcal E\), its maximum property, the given technical context, and the
two displayed numerical hypotheses.  If its first alternative holds, it is
exactly the required elementarity of \(V(T)\), so suppose that its second
alternative holds.  Choose the objects supplied there: a tree-sequence
\(U\), a \noderef{CorrespondingExtendedTashkinovTree} certificate \(C_U\)
relative to the fixed original ETT \(\mathcal E\), maximum property for the
replacement ETT
\[
                   \mathcal E_U=C_U.\mathit{replacementETT},
\]
a generating-color \noderef{HierarchyBase} \(B_U\), a
\noderef{TashkinovHierarchy} \(H\) over that base, and a
\noderef{GoodHierarchy} witness.  We also retain the asserted equality
\[
 V(U)=V(F),\qquad
 F=\mathcal E.\mathit{series.data.tree}
       (\operatorname{Fin.last}(n+1)).                         \tag{1}
\]
The supplied terminal closedness will not be needed in this consumer.

We verify the strict-lower package required by
\noderef{GoodHierarchyElementaryLarge}.  Let \(\mathcal L\) be any
positive-rung maximum ETT, with arbitrary certified input and generating
colorings, whose rung number \(m\) satisfies
\(m<\mathcal E_U.\mathit{rounds}\).  The rung-agreement field in the
corresponding certificate gives
\[
             \mathcal E_U.\mathit{rounds}=n,
\]
so \(0<m<n\).  Apply, respectively, the six fields
\texttt{lowerOne}, \texttt{lowerTwo}, \texttt{lowerThree},
\texttt{lowerFour}, \texttt{lowerFive}, and \texttt{lowerSix} of the
given \noderef{TechnicalInductionContext} to \(\mathcal L\), its positivity,
the transported strict inequality \(m<n\), and its maximum property.  By
the statement of that context, these are exactly assertions (i)--(vi) for
\(\mathcal L\), with all the dependent terminal trees, stored sets,
colorings, membership witnesses, and first-appearance indices belonging to
\(\mathcal L\).  Their conjunction is therefore the full strict-lower
package demanded by the good-hierarchy theorem.

Next let \(\mathcal L\) be any maximum ETT whose rung equals that of
\(\mathcal E_U\).  Composing this equality with the same rung-agreement
field gives
\(\mathcal L.\mathit{rounds}=n\).  The fields
\texttt{currentThree}, \texttt{currentFour}, \texttt{currentFive}, and
\texttt{currentSix} of the technical context, applied to this equality and
the maximum property of \(\mathcal L\), give exactly assertions
(iii)--(vi) at the current rung.  This supplies the second package required
by \noderef{GoodHierarchyElementaryLarge}, uniformly over the arbitrary
input and generating colorings of \(\mathcal L\).

It remains to supply that theorem's candidate-uniform missing-color bound.
Take any such same-rung maximum ETT \(\mathcal L\), together with any
hierarchy base, hierarchy, and good witness offered by the premise of the
good-hierarchy theorem.  Discard those three hierarchy objects.  The rung
equality just proved transports the original technical context from rung
\(n\) to the rung of \(\mathcal L\).  Apply
\noderef{TechnicalCurrentLadderMissingBoundLarge} directly to
\(\mathcal L\), its maximum property, that transported context, and the
unchanged degree and large-order hypotheses.  If \(L_m\) is its exact
terminal ladder tree and \(\psi\) its generating coloring, the conclusion is
\[
  2\mathcal L.\mathit{rounds}+11
       \leq |\noderef{TreeMissingColors}(\psi,L_m)|,             \tag{2}
\]
which is precisely the required candidate-uniform estimate.  Thus (2) was
derived without using the supplied base, hierarchy, or good witness.

All hypotheses of \noderef{GoodHierarchyElementaryLarge} now hold for
\(\mathcal E_U,B_U,H\): its maximum property and good hierarchy came from
the construction alternative, the two induction packages were assembled
above, the large-order hypothesis is the original one, and (2) is uniform
over every same-rung candidate.  Hence
\[
 \noderef{ColoredVertexSetElementary}(G,\varphi,V(U)).           \tag{3}
\]

Finally, the \texttt{belowFinal} field of
\noderef{ExtendedTashkinovTree} says that \(T\) is an exact
\noderef{TreeSequenceInitialSegment} of \(F\).  By
\noderef{TreeSequenceInitialSegmentVertexSetSubset},
\(V(T)\subseteq V(F)\); rewriting with (1) gives
\(V(T)\subseteq V(U)\).  For any distinct \(x,y\in V(T)\), this inclusion
places both vertices in \(V(U)\), and (3) says that their missing-color sets
under \(\varphi\) are disjoint.  The defining pairwise-disjoint condition in
\noderef{ColoredVertexSetElementary} therefore holds on \(V(T)\), as
claimed.
\end{proof}
